\documentclass[12pt]{article} 

\usepackage[top=1.5cm, bottom=1.5cm, left=1.5cm, right=1.5cm]{geometry}
\usepackage{color,graphicx,amsmath,mathtools,amssymb,amsthm}
\pagestyle{plain}

\begin{document}  
\pagestyle{empty}
 




%\vskip0.5in
%
%
%\noindent {\small This document is property of Northeastern University.  Unauthorized distribution of the content is not~permitted.}
%
%\vskip0.5cm





\begin{enumerate}
\item
1. Let G be a plane graph with V vertices, E edges, F faces and C components.
Prove that $$V - E + F - C = 1$$

Induction on C:\\

Base case: If there is there is only 1 component, then we have, by Euler's formula, that $$V - E + F - C = V - E + F - 1 = 1$$
Induction: Assume that V - E + F - C = 1 holds for all planar graph components with fewer than k components. 
Let there be k planar graph components. An edge is then added between two of the components to connect them this decreases the amount of components by 1. 
This edge will add 1 to E and subtract 1 from C, but will not change V or F. Thus, we have that\\ 
$$V - (E + 1) + F - (C - 1) = 1$$
which simplifies to\\
$$V - E + F - C = 1$$\\
Thus, by induction, we have that V - E + F - C = 1 for all planar graphs.\\



\end{enumerate} 




\end{document}